\chapter{Interface utilisateur et front-end}
\label{ch:frontend}

\section{Chaîne de build Vite}

Fichiers d'entrée :
\begin{itemize}
  \item \texttt{resources/css/app.css} — Tailwind v4, variantes \texttt{dark}
  \item \texttt{resources/js/app.js} — importe \texttt{theme.js}, \texttt{navbar.js}, axios
\end{itemize}

\section{Composants Blade réutilisables}

\begin{table}[H]
\centering
\caption{Composants principaux}
\begin{tabular}{@{}lp{8.5cm}@{}}
\toprule
\textbf{Composant} & \textbf{Rôle} \\
\midrule
\texttt{navbar} & Navigation ; variantes \texttt{default} / \texttt{chat} \\
\texttt{chatbot} & FAB + fenêtre chat ; POST AJAX \\
\texttt{layouts/app} & Shell principal + Vite \\
\texttt{layouts/chat} & Fil messaging sombre \\
\texttt{theme-toggle} & Bascule clair/sombre \\
\texttt{alert} & Messages flash session \\
\bottomrule
\end{tabular}
\end{table}

\section{Thème clair / sombre}

Clé \texttt{localStorage} : \texttt{sultan-theme}. Le script \texttt{theme.js} applique la classe \texttt{dark} et l'attribut \texttt{data-theme} sur \texttt{<html>} avant le premier paint via \texttt{theme-init}.

\section{Pages à JavaScript inline volumineux}

Les vues \texttt{hotels/index}, \texttt{map/index} et \texttt{restaurants/index} contiennent CSS et JS intégrés pour autonomie visuelle (design « magazine »). Trade-off : maintenabilité vs maquette riche.

\section{Accessibilite et responsive}

\begin{itemize}
  \item Navbar avec tiroir mobile (\texttt{navbar.js})
  \item Comportement scroll-compact sur la barre de navigation
  \item Grilles adaptatives Tailwind / CSS custom sur les listings
\end{itemize}

\section{Design system (tokens)}

\begin{table}[H]
\centering
\caption{Palette et identite visuelle}
\begin{tabular}{@{}lll@{}}
\toprule
\textbf{Token} & \textbf{Couleur} & \textbf{Usage} \\
\midrule
SultanTeal & \texttt{\#0D4F4F} & Titres, liens, structure \\
SultanGold & \texttt{\#C9A227} & Accents premium \\
SultanTerracotta & \texttt{\#B85C38} & CTA, fleches flux \\
SultanSand & \texttt{\#F5F0E8} & Fonds cartes, boites info \\
\bottomrule
\end{tabular}
\end{table}

\section{Performance front-end}

\begin{itemize}
  \item Build production : \texttt{npm run build} (assets hashes dans \texttt{public/build}).
  \item Pages lourdes : le JS inline augmente le HTML initial ; mitigation future = extraction modules ES.
  \item CDN Leaflet : dependance reseau ; envisager self-host en prod.
\end{itemize}
